\chapter{Cauchy deconvolution}\label{chap:cd}\chaptermark{Cauchy deconvolution}

Every computation in Chapters~\ref{chap:lo} to~\ref{chap:es} inverts the same operator: convolution with the Cauchy density $P_a(x)=(a/\pi)/(a^2+x^2)$. The Thorin measure at level $\theta'$ is the solution, when a positive one exists, of $P_{1-\theta'}*\mu=u_1$; the explicit formula is that equation at $\theta'=\tfrac12$, in measures under RH (Theorem~\ref{tf:thm:deconv}(b)) and in distributions unconditionally (Theorem~\ref{dc:thm:weil}); the matrix pencil solves it for atoms. This chapter collects what is known about the operator in its own right, proves the facts the book uses, and places the book's results in the literature.

The facts fit in a few sentences. Convolution with $P_a$ is harmonic extension to height $a$ in the upper half plane, so Cauchy deconvolution is the Cauchy problem for Laplace's equation, Hadamard's original example of an ill posed problem \cite{Hadamard23}. It is unique when it is solvable. The forward map preserves positivity and the inverse does not: positive data can come from a signed source, and whether a positive source exists is a statement about positive harmonic functions, which Herglotz and Riesz characterised. Positivity of the deconvolution is monotone in the width, so every positive datum has a critical width beyond which no positive source exists; for the Thorin family that width is $1$ minus the supremum of the real parts of the zeros. Stability is H\"older between levels, with the exponent given by the three lines theorem, and logarithmic at the boundary, and both rates are attained. In statistics Cauchy noise is supersmooth of order one, and the minimax rates of density deconvolution are logarithmic; Bayesian posteriors contract at logarithmic rates in the cases treated (Section~\ref{cd:sec:bayes}). All of this changes when the source is a sum of atoms: the problem becomes Prony's, and the error is Lipschitz in the data error, with a constant that grows polynomially as atoms approach, except near merging atoms, where it is the square root.

The chapter develops a survey chapter on Cauchy deconvolution by Nick Polson (the last of the source chapters, of which Chapters~\ref{chap:lv} to~\ref{chap:es} are the earlier ones), together with the checks made for this book. The author's figure is not reproduced: its two panels are described in words and in Table~\ref{cd:tab:slopes}. What is new in the book is the generality of the statements (a datum that is any nonnegative function, a source that is any measure of the class $\int(1+x^2)^{-1}d\lvert\mu\rvert<\infty$), the example of Section~\ref{cd:sec:positivity}, the three lines bound with its sharpness, the bound for finite sources (Proposition~\ref{cd:prop:finite}), the positive two level bound (Corollary~\ref{cd:cor:positive}), the logarithmic bound and the survey of the statistical, algorithmic and physical literature. The specialisation to the Thorin family, the uniqueness argument and the unboundedness of the deconvolution are in Chapters~\ref{chap:tf} and~\ref{chap:dc} and are cited.

The notation is the book's, with the following warnings. A level $\theta$ is the real part of a point $s=\theta+iv$, and $u_\theta(v)=\pi^{-1}\Re(\xi'/\xi)(\theta+iv)$ is the density of Chapters~\ref{chap:lo} to~\ref{chap:es}; in the notation of Chapter~\ref{chap:tf}, which extends equation~\eqref{eq:rho}, it is $\rho_{\theta-1/2}(v)$, the ordinate being $v$ here. The measure $\mathrm{Cauchy}(\gamma,a)$ of Chapter~\ref{chap:lv} is $P_a(\cdot-\gamma)\,dx$. The Thorin measure of Chapters~\ref{chap:lo} to~\ref{chap:es} is $\mathfrak U_\theta=\sum_\rho\mathrm{Cauchy}(\gamma_\rho,\theta-\sigma_\rho)$ with density $u_\theta$, and $\mathfrak U_\theta=2\widetilde U_\theta$ with $\widetilde U_\theta$ the measure of Theorem~\ref{lv:thm:semigroup}, which has the factor $\tfrac12$. The harmonic extension of a measure is written $h$, and not $U$ as in the source. The largest real part of a zero is $\sigma_{\max}=\sup\Re\rho$ (the source's $\Theta$). Widths are $a,c,d,w\ge0$, except that $b$ may be negative with $b>-w$ (Theorem~\ref{cd:thm:holder}, Proposition~\ref{cd:prop:finite}), the example of Section~\ref{cd:sec:positivity} has its own letters $q$, $r$ and $m$, and $\eta$ is the linear coefficient of Theorem~\ref{cd:thm:herglotz}. Throughout, a measure $\mu$ is called admissible when $\int(1+x^2)^{-1}d\lvert\mu\rvert(x)<\infty$, the zero counting measure $\sum_\gamma\delta_\gamma$ being the example that matters.

Status words follow the rest of the book: proved, computed, conditional, heuristic or open. Statements taken from the literature are marked cited and carry the reference; the notes of the last section name the entry that the citation verifier has not confirmed and the statements not checked against the papers. The results are these.
\begin{enumerate}
\item Uniqueness of Cauchy deconvolution, the description of the range of $P_a*$ on $L^2$ and the failure of the na\"ive description of it (Propositions~\ref{cd:prop:unique} and~\ref{cd:prop:range}, proved).
\item A positive datum with a signed source, with the closed form of its critical width (Example~\ref{cd:ex:signed}, proved; the numbers are computed), the critical width of a general positive datum, attained and possibly infinite (Proposition~\ref{cd:prop:critical}, proved), $c^*(u_1)=1-\sigma_{\max}$ (Theorem~\ref{tf:thm:deconv}(a), cited), and the Herglotz characterisation of the data that have a positive source (Theorem~\ref{cd:thm:herglotz}, proved), with Lagarias's criterion as its case $c=\tfrac12$ (cited).
\item H\"older stability between levels, sharp (Theorem~\ref{cd:thm:holder}, proved); the same for finite signed sources with the sharp exponent $(b+w)/(a+w)$ (Proposition~\ref{cd:prop:finite}, proved); a two level bound for positive sources on the line (Corollary~\ref{cd:cor:positive}, proved); the logarithmic bound at the boundary, sharp in order and in the constant (Theorem~\ref{cd:thm:log}, proved). The explanation of why $\widetilde U_{3/4}$ is recovered better than $\widetilde U_{1/2}$ is heuristic, and a test of the exponent on Lorentzians is computed (Table~\ref{cd:tab:slopes}).
\item A survey of density deconvolution, Bayesian inverse problems and positivity preserving algorithms (cited), and a polynomial upper rate for sources analytic in a strip (Proposition~\ref{cd:prop:strip}, proved).
\item Atomic sources and merging atoms, and the analytic continuation problem of many body physics (cited, with the proofs of Chapters~\ref{chap:cx} and~\ref{chap:es}), and three open problems.
\end{enumerate}
